\section{Phase comparisons and new forcing regions}
\label{sec:phase-comparisons}
We work with the equation
\begin{equation}\label{ph:eq}
 X^iY^2X^kYX^\ell=U^iV^2U^kVU^\ell,
 \qquad k\ge1,\quad i+\ell\ge1.
\end{equation}
No maximum-length condition is imposed. The result below gives direct
proofs in infinite exponent regions, rather than a bound on the length
of a counterexample. The remaining exponent triples are not settled here.

\begin{theorem}\label{ph:main}
Equation \eqref{ph:eq} forces periodicity in each of the following cases:
\begin{enumerate}
\item $\ell=0$, $i,k\ge1$, and $k\ne2i$;
\item $i=0$, $\ell\ge1$, and $k>2\ell$;
\item $i,k,\ell\ge1$, $2k>i+\ell$, and $i+k>2\ell$;
\item $i,k,\ell\ge1$, $k\ge2$, $i+k>2\ell$, and $2k+2\ell>i$.
\end{enumerate}
The only exception on the first slice is $k=2i$, for which family D in
Theorem~\ref{nonforcing} supplies a counterexample satisfying Holub's
length condition. In particular the whole slice $k>2(i+\ell)$ with
$i+\ell>0$ is forcing, even without that condition.
\end{theorem}

\subsection{Preliminaries and notation}
We may exchange the two substitutions and assume $d=x-u>0$;
otherwise equal lengths give identical substitutions. Put
$S=i+k+\ell$, $D=v-y=Sd/3>0$. By Lemma~\ref{perred}, none of
these tuples can have a counterexample with a periodic morphism.
In particular all four images are nonempty in a counterexample. (Every
empty image makes its morphism periodic.) Equivalently, the same
reduction follows from the Spehner binary-code power theorem: the
pattern is primitive, has three $b$'s and at least two $a$'s, and the
image under a nonperiodic morphism cannot be a proper power.
The exceptional seed $i=k=1,\ell=0$ is already covered by
Theorem~\ref{rays}.

Let $X=R^t$ with $R$ primitive, $r=|R|>0$, $t\ge1$. Index the
bi-infinite periodic word $R^{\mathbb Z}$ so that position zero begins a
copy of $R$. By the \emph{phase} of a factor we mean its starting or
ending position modulo $r$, as specified. Two equal factors of length
at least $r$ have the same starting phase (and the same ending phase),
since the $r$ rotations of a primitive word are distinct. We will
compare factors of separate powers of $X$ in this coordinate system.

We use two classical facts. First, if a factor of $R^{\mathbb Z}$ equals
$U^m$ and has length at least $r+u-\gcd(r,u)$, then Fine--Wilf gives
period $\gcd(r,u)$. Its length is at least $r$ and $u$, so the primitive
roots of $U$ and $R$ are conjugate and $r\mid u$. If $U$ is also a prefix
or suffix of $X$, those roots are equal, hence $U\in R^*$. This prefix
or suffix property holds in \eqref{ph:eq} as $i+\ell>0$ and $x>u$.
Second, two noncommuting words freely generate a rank-two subgroup of
the ambient free group (see Hadravov\'a--Holub, \cite{hh}, Lemma 1).
Thus, if $R,U$ do not commute, on $\langle R,U\rangle$ there is a
homomorphism $\nu$ to $\mathbb Z$ with $\nu(R)=0$, $\nu(U)=1$.
We use this exponent-count homomorphism only after expressing all
relevant words as group products of $R$ and $U$. All original word
equalities remain literal monoid equalities; inverses are used only in
these auxiliary group computations.

\paragraph{Transport of a prefix or suffix.}
If $B$ is a positive power of $R$ and $Y$ is a prefix of $BY$, then $Y$
is a prefix of $R^\omega$, and $BY=YZ$ where $Z$ is a power of the
rotation of $R$ at phase $y$. The prefix assertion follows by repeatedly
using $Y[j]=Y[j-|B|]$ for $j\ge |B|$; the equation and rotation follow
by reading $BY$ on $R^\omega$. Dually, if $Y$ is a suffix of $YB$,
then $Y$ is a suffix of a sufficiently long power of $R$, and $YB=ZY$
where $Z$ is a power of the rotation of $R$ at phase $-y$.

\subsection{An exact cancellation identity in the nested region}
Put
\[
 L=2D-id,\quad Q=D-\ell d,\quad
 H=Q-id=\frac{k-2(i+\ell)}3d,\quad
 F=L-\ell d=\frac{2k-i-\ell}3d.
\]
Here $F-H=D$. If $L,Q\ge0$, the middle $U^k$ lies wholly inside the
middle $X^k$ (the left minus right boundary displacements are
$id,id-D,id-2D,(i+k)d-2D,-\ell d,0$). Define words $C,B,P,A$ by
\begin{equation}\label{ph:factor}
 X^i=U^i C,\qquad X^\ell=B U^\ell,\qquad
 X^k=P U^k A.
\end{equation}
Their lengths are $id,\ell d,L,Q$, respectively. For zero outer
exponents the corresponding word is empty. Comparing suffixes of the
common expansion and cancelling $U^\ell$ gives $V=AYB$. Cancelling the
prefix $U^i$ and the common suffix $U^k AYB U^\ell$ gives
\begin{equation}\label{ph:cancel}
 C Y^2 P=(A Y B)^2.
\end{equation}
In particular $C,A$ are prefix-comparable and $P,B$ are suffix-comparable.
These comparisons are of literal words, independent of the size of $Y$.

\begin{lemma}\label{ph:phase}
In the nested region, if $i>0$ and
\[
 \min(iu,ku)+\min(id,Q)\ge r,
\]
then $r\mid H$. If $\ell>0$ and
\[
 \min(\ell u,ku)+\min(\ell d,L)\ge r,
\]
then $r\mid F$.
\end{lemma}
\begin{proof}
For the first comparison extend the matching prefixes of $C,A$ on
the left by the common suffix of $U^i,U^k$, of length $\min(iu,ku)$.
The resulting two equal factors of powers of $X$ have phases that
differ by $(-Q)-(-id)=-H$. Their length is at least $r$. For the
second, extend the matching suffixes of $P,B$ on the right by the
common prefix of $U^k,U^\ell$. Their ending phases differ by $L-\ell
d=F$. Apply primitive rotation uniqueness.
\end{proof}

\begin{lemma}\label{ph:prop}
In the nested region, either $r\mid D$ or $XU=UX$ suffices to force
periodicity.
\end{lemma}
\begin{proof}
If $r\mid D$, then $r\mid Sd$ and hence $r\mid Su$ since $r\mid x$.
The factor $U^i$ starts at phase zero, $U^k$ at phase $L\equiv iu$
(modulo $r$), and $U^\ell$ starts at phase $-\ell u\equiv(i+k)u$.
For zero exponents the corresponding condition is vacuous. Therefore
$U^S$ equals the prefix of $R^\omega$ of length $Su$, which is a
positive power of $R$. Uniqueness of primitive roots gives $U\in R^*$.
Conversely, if $XU=UX$, then $r\mid u,d$ and the occurrence of the
full word $U^k$ at phase $L$ yields $r\mid L$. So $r\mid2D$ and
$r\mid3D=Sd$, whence $r\mid D$.

In either case $r$ divides $Q$ and $\ell d$. Thus
$V=AYB=R^{Q/r}Y R^{\ell d/r}$. If $R,Y$ do not commute they form a
two-word code, and substituting this expression into \eqref{ph:eq}
gives two unequal sequences of codewords: the left has adjacent
$Y$ tokens while the first two $Y$ tokens on the right are separated
by $(Q+\ell d)/r=D/r>0$ copies of $R$. Injectivity is contradicted.
Consequently $Y,V\in R^*$ too.
\end{proof}

\subsection{Positive outer exponents}
Suppose $i,k,\ell\ge1$. First assume $Q,L\ge d$. Both lengths in
Lemma~\ref{ph:phase} are at least $u+d=x\ge r$. Hence $r\mid H,F$,
so $r\mid F-H=D$, and Lemma~\ref{ph:prop} applies. This proves a useful subregion of part 4 of Theorem~\ref{ph:main}.

For part 3, put $c=i+\ell$, $M=2k-c>0$. We have $Q>0$ and
$L=F+\ell d>\ell d$, so the nested identity applies. The second
comparison in Lemma~\ref{ph:phase} has length at least $u+d=x$,
whence $r\mid F=Md/3$. In particular $r\mid Md$, and $r\mid Mu$.
If $M\le k-1$, Fine--Wilf on $U^k$ forces $U\in R^*$, since
$ku\ge r+u$. If $M=k$ and $r<ku$, then $r$ is a proper divisor
of $ku$, so $r\le ku/2\le(k-1)u$ ($k\ge2$). If $M=k$ and $r=ku$,
the length-$r$ factor $U^k$ would be an imprimitive rotation of $R$,
a contradiction. Thus all cases $M\le k$ are settled.

It remains to consider $M>k$ (equivalently $k>c$) and assume
$U\notin R^*$. The same Fine--Wilf argument gives $r>(k-1)u$.
Since $c\ge2$, $M\le2(k-1)$; a proper divisor of $Mu$ is at most
$(k-1)u$. Consequently $r=Mu$. Since $r\mid F$, the start of
$U^k$ has phase $L\equiv\ell d\equiv-\ell u$. As $k>\ell$ and
$ku<r$, the portion crossing phase zero consists of $U^\ell$ at the
end and $U^{k-\ell}$ at the beginning of $R$. In particular
\[
 R=U^{k-\ell}G U^\ell,
 \qquad |G|=(M-k)u=(k-c)u>0.
\]
Here $A$ begins at phase $(k-\ell)u$, so begins with $G$ if its
length permits; $C$ begins at phase $iu<(k-\ell)u$, so begins with
$U^{k-c}$ if its length permits. We check both lengths explicitly.
Since $x\ge r$, $d\ge(M-1)u$, and $id\ge(k-c)u$. If $k\ge2c$,
then $Q\ge id$ and hence $Q\ge(k-c)u$. If $k<2c$, then
$Q\ge d/3\ge(M-1)u/3\ge(k-c)u$, the last inequality being
$k\le2c-1$. Prefix comparability in \eqref{ph:cancel} now gives
$G=U^{k-c}$, contradicting the primitivity of $R$ ($M>1$).
This proves part 3.

\subsection{The rest of the strictly nested region with $k\ge2$}
We prove part 4. Thus $Q,L>0$ and $L+Q=kd\ge2d$; at least one of
$Q,L$ is at least $d$. The case $F>0$ has just been proved in part 3.
If $F=0$, then $i+\ell=2k$ and $S=3k$. The inequality $Q>0$ implies
$i>\ell$, so $k<i<2k$ and $i+k>2k$. The balanced-prefix theorem
(Theorem~\ref{balanced}) applies, and none of its nonforcing exceptions
is in this region. Hence we assume $F<0$.

If $L\ge d$, the second comparison in Lemma~\ref{ph:phase} gives
$r\mid F$. Put $M=i+\ell-2k>0$, $m=\max(i,k,\ell)\ge2$. Then
$r\mid Mu$. A run of $U^m$ is a factor of a power of $X$ (the outer
runs are common prefixes or suffixes; the middle is nested).
If $r\le(m-1)u$, Fine--Wilf settles the case. If $M<m$, this is
immediate; if $M=m$ and $r=mu$, $U^m$ is a length-$r$ factor of
$R^{\mathbb Z}$, impossible by primitivity. Otherwise $M>m$ and
$M\le2m-2$ because $k\ge1$. The only divisor of $Mu$ greater than
$(m-1)u$ is $Mu$ itself. Thus $r=Mu$. Both $i,\ell<M$, and $R$
has prefix $U^i$ and suffix $U^\ell$. These cover $R$, since
$i+\ell=M+2k>M$, and their $u$-letter blocks are aligned because
$r=Mu$. This makes $R=U^M$, impossible. We have shown that $r\mid F$
with $F<0$ always suffices in the nested region.

It remains that $Q\ge d$ (in particular if $L<d$). The first
comparison gives $r\mid H$. Write $C=AE$, $B=GP$ using comparability.
Here $|E|=-H=D-F>D$, $|G|=-F>0$, and $E=R^e$ with $e\ge1$.
Cancel $A$ and the terminal $P$ in \eqref{ph:cancel} to obtain
\begin{equation}\label{ph:negative}
 E Y^2=Y B A Y G.
\end{equation}
Transport $EY=YZ$ with $Y$ a prefix of $R^\omega$, $Z$ a power
of the rotation at phase $y$. Then $ZY=B A Y G$.
In particular $T=BA$, of length $D<|Z|$, is a prefix of $Z$.
If $\ell d\ge r$, the occurrence of $B$ at the start of $T$ and
at phase zero in $X^\ell$ gives $r\mid y$. Hence $Y\in R^*$.
As $r\mid H$, the phase of $A$ is $-Q\equiv iu$, so
$A=U^{-i}R^a$ in the free group. From \eqref{ph:factor} we have
$B=R^{t\ell}U^{-\ell}$ and $P=R^{tk-a}U^{i-k}$. Hence
$G=B P^{-1}=R^{t\ell}U^{k-i-\ell}R^{a-tk}$.
If $R,U$ do not commute, applying $\nu$ to \eqref{ph:negative}
yields $0=k-2(i+\ell)$, impossible for $F<0$. Otherwise apply
Lemma~\ref{ph:prop}.

Suppose finally that $\ell d<r$ and $R,U$ do not commute. Since
$Q\ge d$, we have $i+k\ge2\ell+3$. If $\ell\ge2$, the largest
of $i,k$ is at least four. Fine--Wilf on that run then forces
$u<d/2$, so $r\le x<3d/2<\ell d$, a contradiction. Thus $\ell=1$.
The conditions $F<0$ and $Q\ge d$ give $i\ge2k$ and $i+k\ge5$,
so $i\ge4$. Fine--Wilf on $U^i$ gives both $r>(i-1)u$ and
$(i-2)u<d$, so $r\le x<2d$; as $r>d$ and $r\mid x$, $X=R$.
The divisor condition $r\mid H$ implies $r\mid M'u$ where
$M'=2i+2-k$. Since $M'\le2i+1\le3(i-1)$ for $i\ge4$,
the quotient $M'u/r$ is one or two. If two, then
$r=(i+1-k/2)u$. For $k\ge4$, this is at most $(i-1)u$,
impossible. For $k=3$ the Fine--Wilf threshold is $iu$ exactly
($\gcd(r,u)=u/2$), and for $k=2$ the prefix $U^i$ equals $R$,
impossible. For $k=1$ (not needed for part 4, but included in
this argument), $r=(i+1/2)u$ and the prefix $U^i$ overlaps the
suffix $U$ by $u/2$. The two halves of $U$ are equal, so the
prefix and suffix make $R$ a power of this half, again impossible.

If the quotient is one, $r=M'u$, and $R$ has prefix $U^i$ and
suffix $U$, so $R=U^i K U$ with $|K|=(i+1-k)u=gu$ say. Here
$1\le g\le i$. The phase of $A$ is $iu$, hence $A$ begins with
$KU$ (it is a suffix of powers of $R$ with that phase).
Also $B$ is $R$ with its final $U$ removed. The factor $T=BA$
of $R^{\mathbb Z}$ therefore begins with $U^i K K U$, of length
$r+gu$. Its period $r$ makes its first $gu$ letters, namely
$U^g$, equal its letters at positions $r$ through $r+gu-1$,
namely the suffix of $K$ after its first $u$ letters followed by
$U$. Consequently $K=J U^{g-1}$ with $|J|=u$.
The first $r$ letters of $T$ now have the form $U^i J U^{g-1}J$,
a cyclic rotation of $R=U^i J U^g$. There are $M'=i+g+1>4$
blocks of length $u$. If $J\ne U$, a rotation of $R$ has at least
$M'-2$ blocks equal to the same rotation of $U$, since at most
two blocks touch $J$. The proposed rotation has exactly two $J$
blocks and $M'-2>2$ blocks equal to $U$. The repeated rotated
block must be $U$, so the rotation fixes $U^{M'}$. Its Hamming
distance from $U^{M'}$ must then be invariant, but it changes
from $h$ to $2h$, where $h>0$ is the distance between $J,U$.
Thus $J=U$ and $R$ is imprimitive, a contradiction.
This completes part 4.

\subsection{The zero-tail slice $\ell=0$}
Assume $i,k\ge1$, $k\ne2i$, and $i+k\ge3$; the excluded small seed
was already proved. Then $D\ge d$. First suppose $i\le2k$, so
$L\ge0$, $Q=D>0$ and the nested factorization applies. Here $B$ is
empty. Lemma~\ref{ph:phase} gives $r\mid H=D-id$, since its first
comparison has length at least $u+d$. Consequently
\begin{equation}\label{ph:zeroarithmetic}
 r\mid (k-2i)d,\qquad r\mid (k-2i)u.
\end{equation}
If $k>i$, then $0<|k-2i|\le k-1$, giving $r\le(k-1)u$ and
Fine--Wilf on $U^k$ followed by Lemma~\ref{ph:prop} settles the case.

Now $k\le i\le2k$. Set $C=A E$ (prefix comparison), so
$|E|=id-D>0$ is a multiple of $r$. Since $C$ ends at phase zero,
$E=R^e$, $e\ge1$. Cancelling $A$ in \eqref{ph:cancel} gives
\begin{equation}\label{ph:tail1}
 E Y^2 P=Y A Y.
\end{equation}
Then $Y$ is a prefix of $EY$; transport gives $EY=YZ$ with $Z$
of phase $y$ and length $|E|$. Cancelling $Y$ shows $ZYP=AY$.
As $|E|=id-D\le D=|A|$, $Z$ is a prefix of $A$. Comparing at
least $r$ letters gives $y\equiv-D\equiv iu\pmod r$.

Both $Y$ and $U^i$ are prefixes of $R^\omega$ with lengths congruent
modulo $r$, and $A$ is a suffix of a power of $R$ of length $D$.
In the free group there exist integers $b,a$ such that
\[
 Y=R^b U^i,\qquad A=U^{-i}R^a.
\]
Indeed $U^i A$ reads continuously on $R^\omega$ and has length a
multiple of $r$. Moreover $X^k=P U^k A$ implies
$P=R^{tk-a}U^{i-k}$. If $R,U$ do not commute, apply $\nu$ to
\eqref{ph:tail1}: $2i+(i-k)=i-i+i$, so $k=2i$, excluded.
Thus they commute. Lemma~\ref{ph:prop} concludes this nested case.

For $i>2k$, the middle run is not nested. Define $C$ by
$X^i=U^i C$ and write $V=AY$ from the final suffix comparison
($Y$ is a suffix of $V$), $|A|=D$. Then $C=A E$ with
$|E|=id-D>D>0$. Let $J$ be the left overflow just after the second
$Y$ (length $id-2D>0$). Corresponding prefix cancellations give
\begin{equation}\label{ph:tail2}
 E Y^2=Y A Y J,\qquad J X^k=U^k A.
\end{equation}
These also follow by cancelling the common $U^i$ and final $Y$
and comparing the intermediate prefix lengths.

We claim $r\mid |E|$. If $D\ge kx$, then $X^k$ is a suffix of $A$
by \eqref{ph:tail2}; this full power occurs in $X^i$ at phase
$iu+D-kx$, yielding $r\mid iu+D$, hence $r\mid id-D$. If $D<kx$,
then $A$ is a suffix of $X^k$. The word before $A$ in $X^k$ is a
suffix of $U^k$, of length $kx-D$, since $J X^k=U^k A$ and $J$
is longer than zero. Compare the two occurrences of $A$ in $X^i$
and $X^k$, extending to the left by $\min(iu,kx-D)$ common
letters from $U^i,U^k$; the total length is
$\min(D+iu,kx)\ge x$ because $D\ge d$. The phases yield the same
conclusion. In either case $E=R^e$ with $e\ge1$, and $A$ is a
suffix of a power of $R$ (from $X^i=U^i A E$).

We next show that $D\ge r$ unless $R,U$ already commute. If they do
not commute, the common prefix $U^i$ cannot satisfy the Fine--Wilf
threshold with periods $r,u$, so $(i-2)u<d$ because $r\le x=u+d$.
If $i\ge4$, $x<3d/2$ while $D\ge5d/3$, hence $D>r$.
Otherwise $i>2k$ implies $(i,k)=(3,1)$, so $D=4d/3$, $|E|=5d/3$.
If $D<r$, divisibility gives $r=5d/3$. Since $x<2d$ and $r\mid x$,
$x=r$, $u=2d/3$. Then $3u=r+u-\gcd(r,u)$, the Fine--Wilf
threshold, contradicting noncommutation. If $R,U$ already commute,
$r\mid d$ and $D\ge d\ge r$. Thus in every case $D\ge r$.

From \eqref{ph:tail2}, $Y$ is a prefix of $EY$; again transport
$EY=YZ$ with phase $y$. The equality $ZY=AYJ$ and $|Z|>|A|\ge r$
give $y\equiv-D\equiv iu\pmod r$. As before in the free group
$Y=R^b U^i$, $A=U^{-i}R^a$, and now
$J=U^{k-i}R^{a-tk}$. If $R,U$ do not commute, applying $\nu$ to
$E Y^2=Y A Y J$ yields $2i=i-i+i+k-i=k$, impossible. Thus
$R,U$ commute. Now $r\mid d$ and $r\mid id-D$ give $r\mid D$;
$y\equiv-D$ yields $r\mid y$ and $Y\in R^*$, and $V=AY\in R^*$.
This finishes the whole zero-tail slice.

\subsection{The zero-head slice $i=0$, $k>2\ell>0$}
Now $L=2D$, $Q=(k-2\ell)d/3>0$, $F=(2k-\ell)d/3>0$.
The second comparison in Lemma~\ref{ph:phase} yields $r\mid F$,
as $L>\ell d$. By suffix comparability write $P=E B$,
$E=R^e$, $e\ge1$. Note $|E|=F=D+Q$. Equation \eqref{ph:cancel}
reduces after cancelling $B$ to
\begin{equation}\label{ph:head}
 Y^2 E=A Y B A Y.
\end{equation}
Thus $Y$ is a suffix of $YE$. The suffix transport gives
$YE=ZY$, where $Z$ is a rotation of $R^e$ of phase $-y$.
Cancelling the final $Y$ in \eqref{ph:head}, we obtain
$YZ=A Y T$, $T=BA$. Since $|T|=D<|Z|$, $T$ is a suffix of $Z$
and is a factor of $R^{\mathbb Z}$ ending at phase $-y$.

If $Q=|A|\ge r$, compare $A$ as the suffix of $Z$ and the suffix of
$X^k$. This gives $r\mid y$, so $Y,Z$ are powers of $R$.
Now $A$ is a prefix of $YZ$ and has length at least $r$, giving
$r\mid Q$. Since $r\mid F=D+Q$, we have $r\mid D$, and apply
Lemma~\ref{ph:prop}.

If $|B|=\ell d\ge r$, compare $B$ as a prefix of $X^\ell$ and as the
prefix of $T$ in $Z$ at phase $-y-D$. It gives $y\equiv-D$.
Since $r\mid D+Q$, the suffixes $A$ and $Y$ of powers of $R$ have
congruent lengths. Therefore $A=Y R^a$ for some integer $a$ in the
free group. From $B U^\ell=X^\ell$ and $X^k=E B U^k A$, we also
obtain $B=R^{t\ell}U^{-\ell}$ and $A=U^{\ell-k}R^{tk-e-t\ell}$.
In particular $Y$ belongs to $\langle R,U\rangle$. If $R,U$ do not
commute, apply $\nu$ to \eqref{ph:head} to get
$2(\ell-k)=4(\ell-k)-\ell$, hence $\ell=2k$, impossible. If they
commute, apply Lemma~\ref{ph:prop} directly.

It remains to consider $\ell d<r$ and $Q<r$ with $R,U$ not commuting.
Fine--Wilf in the nested middle run gives
$ku<r+u-\gcd(r,u)$; hence $(k-2)u<d$. If $\ell\ge2$ then $k\ge5$
and $x<3d/2<\ell d$, contradiction. Thus $\ell=1$, $k\ge3$.
Now $d<r\le x<2d$, so $X=R$ (a proper divisor of $x$ is at most
$x/2<d$). Because $r\mid F=(2k-1)d/3$ and $r\mid x$, we have
$r\mid(2k-1)u$. Since $r>(k-1)u$, the quotient $(2k-1)u/r$ is
one or two. If two, $r=(k-1/2)u$ with $u$ even and
$\gcd(r,u)=u/2$, giving $ku=r+u-\gcd(r,u)$, a contradiction.
If one, $r=(2k-1)u$, $d=(2k-2)u$ and $F/r=(2k-2)/3$ is integral.
So $k\equiv1$ modulo three. For $k\ge7$,
$Q=2(k-1)(k-2)u/3\ge(2k-1)u=r$, contradiction. Only $k=4$
remains.

Here $r=x=7u$, $d=6u$, $D=10u$, $Q=4u$, $F=14u$.
The middle $U^4$ starts at phase $-u$, so $R$ begins with $U^3$
and ends with $U$. Since $T=BA$ has period $7u$ and length $10u$,
its first and last $3u$ letters agree. Its first $3u$ letters are
$U^3$ (a prefix of $B$), its last $3u$ letters are a suffix of $R$.
Thus $R=U^3 K U^3$ with $|K|=u$. The length-$7u$ prefix of
$T=U^3 K U^2 K U^3$ must be a cyclic rotation of $R$.

For completeness, if $U\ne K$ this is impossible: a rotation of
$U^3 K U^3$ read in seven blocks of size $u$ has at least five blocks
that are the same rotation of $U$, since at most two blocks touch
the unique $K$ block (including cyclic wraparound). But the proposed
rotation $U^3 K U^2 K$ has exactly five $U$ blocks and two $K$
blocks. The repeated rotated block must equal $U$, so this rotation
fixes $U^7$. Cyclic Hamming distance to $U^7$ is invariant under it.
The two words have distances $h$ and $2h$, respectively, where
$h$ is the positive Hamming distance between $K$ and $U$. This is a
contradiction. Consequently $K=U$, and $R=U^7$ is imprimitive.
The final exceptional length regime is thereby eliminated. This
completes the proof of Theorem~\ref{ph:main}.

\paragraph{Status.}
Intersecting Theorem~\ref{ph:main} with the previously unresolved
lists removes $27$ of $63$ tuples through exponent four, and $90$
of $210$ through exponent six. The remaining counts are $36$ and
$120$, respectively. These are unconditional proofs for those removed
tuples, subject to ordinary independent mathematical review; no bound
on $d$ was needed. The other residual tuples, and a polynomial bound
on $d$ for them, remain open in this work.
